\section{讨论}\label{sec:discussion}

这些证据与配置依赖的可蒸馏性相一致，即应相对于特定学生状态、CE 对照及训练协议进行评价。教师模型质量强烈影响迁移，但留出集上的平均值不能完整描述哪些预测目标有助于学生模型。因此，教师状态改变带来的迁移符号反转与较弱监督来源的反例可以同时成立，而不意味着已经得到完整的因果分解。

预测多样性不一定是可迁移的多样性：在存在混杂因素的对照中，分支分歧更大的 TG 教师模型比 TT 提供的迁移收益更小。同样，局部梯度概括的是一个小邻域，而完整训练会通过学习率调度及裁剪设置改变参数。这些不匹配提示，应谨慎将所测试的统计量用作长程预测指标，但并不意味着 KD 结果在一般情况下不可预测。对于已训练的 LM，匹配的 CE 续训仍是相关的操作性对照。

这些对照约束的是不同解释，并非估计可以相互替代的效应。教师状态比较固定了部署时的架构与组成，但改变了分支状态及准备过程。监督来源消融同时改变了组成与似然。TT 对照检验同构替代方案能否达到或超过 TG，而不是一个参数量与质量均匹配的架构对照。共享的 CE 与 TG 最终结果使同一学生内的比较保持一致，但不会成倍增加独立证据。对这些对照的联合解释比任何单独的机制主张都更受限。

检查点候选资格也会改变实际问题。TinyStories 的 CE 续训相较于 $S_0$ 有所改善，因此其负增益是相对于有用续训而言的。在现代模型压力测试中，允许第 0 步参与选择后，每个训练臂均保留 $S_0$；KD 劣于 CE 的现象，与两种训练方式均缺乏进一步续训必要性的情况同时存在。报告这一区别，可避免将所有负增益都视为同一种训练设置的证据。上述对照均不支持用于教师选择或 KD 强度调节的普遍规则。
